\begin{lemma}
\label{lemma-depth-dualizing-module}
Let $(A, \mathfrak m, \kappa)$ be a Noetherian local ring with
normalized dualizing complex $\omega_A^\bullet$. Let $d = \dim(A)$
and $\omega_A = H^{-d}(\omega_A^\bullet)$. Then
\begin{enumerate}
\item the support of $\omega_A$ is the union of the irreducible components
of $\Spec(A)$ of dimension $d$,
\item $\omega_A$ satisfies $(S_2)$, see
Algebra, Definition \ref{algebra-definition-conditions}.
\end{enumerate}
\end{lemma}

\begin{proof}
We will use Lemma \ref{lemma-sitting-in-degrees} without further mention.
By Lemma \ref{lemma-nonvanishing-generically-local} the support
of $\omega_A$ contains the irreducible components of dimension $d$.
Let $\mathfrak p \subset A$ be a prime. By Lemma \ref{lemma-dimension-function}
the complex $(\omega_A^\bullet)_{\mathfrak p}[-\dim(A/\mathfrak p)]$
is a normalized dualizing complex for $A_\mathfrak p$. Hence if
$\dim(A/\mathfrak p) + \dim(A_\mathfrak p) < d$, then
$(\omega_A)_\mathfrak p = 0$.
This proves the support of $\omega_A$ is the union of the irreducible
components of dimension $d$, because the complement of this union
is exactly the primes $\mathfrak p$ of $A$ for which
$\dim(A/\mathfrak p) + \dim(A_\mathfrak p) < d$ as $A$ is catenary
(Lemma \ref{lemma-universally-catenary}).
On the other hand, if $\dim(A/\mathfrak p) + \dim(A_\mathfrak p) = d$, then
$$
(\omega_A)_\mathfrak p =
H^{-\dim(A_\mathfrak p)}\left(
(\omega_A^\bullet)_{\mathfrak p}[-\dim(A/\mathfrak p)] \right)
$$
Hence in order to prove $\omega_A$ has $(S_2)$ it suffices to show that
the depth of $\omega_A$ is at least $\min(\dim(A), 2)$.
We prove this by induction on $\dim(A)$. The case $\dim(A) = 0$ is
trivial.

\medskip\noindent
Assume $\text{depth}(A) > 0$. Choose a nonzerodivisor $f \in \mathfrak m$
and set $B = A/fA$. Then $\dim(B) = \dim(A) - 1$ and we may apply the
induction hypothesis to $B$. By Lemma \ref{lemma-divide-by-nonzerodivisor}
we see that multiplication by $f$ is injective on $\omega_A$ and we get
$\omega_A/f\omega_A \subset \omega_B$. This proves the depth of $\omega_A$
is at least $1$. If $\dim(A) > 1$, then $\dim(B) > 0$ and $\omega_B$
has depth $ > 0$. Hence $\omega_A$ has depth $> 1$ and we conclude in
this case.

\medskip\noindent
Assume $\dim(A) > 0$ and $\text{depth}(A) = 0$. Let
$I = A[\mathfrak m^\infty]$ and set $B = A/I$. Then $B$ has
depth $\geq 1$ and $\omega_A = \omega_B$ by
Lemma \ref{lemma-divide-by-finite-length-ideal}.
Since we proved the result for $\omega_B$ above the proof is done.
\end{proof}